

%%%%%%%%%%%%%%%
\begin{frame}{Requêtes sur la base ``Messagerie''}

Essayons de créer et d'interroger cette instance.

\vfill

\figSlideWithSize{instance-messagerie.pdf}{9cm}

\vfill


Connaissez-vous bien le schéma\,?

\end{frame}


%%%%%%%%%%%%%%%
\begin{frame}[fragile]{Insérons le premier message}

Les contacts:

\begin{minted}[baselinestretch=.9,
               fontsize=\small,
               framesep=2mm]{sql}
insert into Contact (idContact, prénom, nom,  email)
  values (1, 'Serge', 'A.', 'serge.a@inria.fr');
insert into Contact (idContact, prénom, nom,  email)
  values (4, 'Philippe', 'R.', 'philippe.r@cnam.fr');
\end{minted}

\vfill

Le message:

\begin{minted}[baselinestretch=.9,
               fontsize=\small,
               framesep=2mm]{sql}
insert into Message (idMessage, contenu, idEmetteur)
values (1, 'Hello Serge', 4);
\end{minted}

\vfill

L'envoi

\begin{minted}[baselinestretch=.9,
               fontsize=\small,
               framesep=2mm]{sql}
insert into Envoi (idMessage, idDestinataire) values (1, 1);
\end{minted}

\vfill
\begin{blocgauche}
Saurez-vous insérer les autres messages?
\end{blocgauche}
\end{frame}

%%%%%%%%%%%%%%%
\begin{frame}[fragile]{Insérons le second message}

Le message:

\begin{minted}[baselinestretch=.9,
               fontsize=\small,
               framesep=2mm]{sql}
insert into Message (idMessage, contenu, idEmetteur, idPrédecesseur)
values (2, 'Coucou Philippe', 1, 1);
\end{minted}

\vfill

L'envoi

\vfill

\begin{minted}[baselinestretch=.9,
               fontsize=\small,
               framesep=2mm]{sql}
insert into Envoi (idMessage, idDestinataire) 
values (2, 4);
\end{minted}

\vfill
Pas facile de s'y retrouver avec les identifiants... 

\vfill

\alert{Une bonne
interface graphique aiderait beaucoup}.
\end{frame}


%%%%%%%%%%%%%%%
\begin{frame}[fragile]{Interrogeons}

Les messages et leur émetteur.

\begin{minted}[baselinestretch=.9,
               fontsize=\small,
               framesep=2mm]{sql}
select idMessage, contenu, prénom, nom
from Message as m,  Contact as c
where m.idEmetteur = c.idContact
\end{minted}

\vfill

On obtient:

\vfill

\begin{small}
\begin{tabular}{c|c|c|c}
\textbf{idMessage} & \textbf{contenu} & \textbf{prénom} & \textbf{nom}\\ \hline 
1 & Hello Serge & Philippe & R.\\ \hline 
2 & Coucou Philippe & Serge & A.\\ \hline 
\end{tabular}
\end{small}
\end{frame}


%%%%%%%%%%%%%%%
\begin{frame}[fragile]{Interrogeons encore}

Les messages et leur prédecesseur.

\begin{minted}[baselinestretch=.9,
               fontsize=\small,
               framesep=2mm]{sql}
select m1.contenu as 'Contenu', m2.contenu as 'Prédecesseur'
from Message as m1,  Message as m2
where m1.idPrédecesseur = m2.idMessage
\end{minted}

\vfill
\begin{small}
\begin{tabular}{c|c}
\textbf{Contenu} & \textbf{Prédecesseur}\\ \hline 
Coucou Philippe & Hello Serge\\ \hline 
Philippe a dit ... & Hello Serge\\ \hline 
Serge a dit ... & Coucou Philippe\\ \hline 
\end{tabular}
\end{small}
\vfill

Sauriez-vous obtenir \alert{tous} les prédecesseurs
d'un message?
\end{frame}


%%%%%%%%%%%%%%%
\begin{frame}[fragile]{Interrogeons encore et encore}

Les messages envoyés à plus d'un contact

\begin{minted}[baselinestretch=.9,
               fontsize=\small,
               framesep=2mm]{sql}
select m.idMessage, contenu, count(*) as 'nbEnvois'
from Message as m, Envoi as e
where m.idMessage = e.idMessage
group by idMessage, contenu
having nbEnvois > 1
\end{minted}

\vfill

\begin{tabular}{c|c|c}
\textbf{idMessage} & \textbf{contenu} & \textbf{nbEnvois}\\ \hline 
4 & Serge a dit ... & 2\\ \hline 
\end{tabular}

\vfill

Requête effectuée souvent? On pourrait créer une vue.
\end{frame}


%%%%%%%%%%%%%%%
\begin{frame}[fragile]{Nettoyons un peu la base}

Supprimons les vieux messages

\begin{minted}[baselinestretch=.9,
               fontsize=\small,
               framesep=2mm]{sql}
begin transaction;

delete from Message; where year(dateEnvoi) < 2015

Réponse: All messages deleted. Table message is now empty...
\end{minted}


\vfill
Oups... un point-virgule mal placé. \alert{Que faut-il faire\,?}

\vfill

Partir très loin\,? Arracher tous les fils\? ou ... ?

\end{frame}
%%%%%%%%%%
\begin{frame}{À retenir}

Les requêtes SQL ``en direct'' sont assez peu pratiquées.

\vfill

\begin{redbullet}
  \item La manipulation des identifiants est peu commode
  \item Risque important de faire une bêtise
  \item \alert{Utile pour tester des requêtes, ou vérifier la base}
\end{redbullet}

\vfill

En général, les requêtes sont intégrées à des applications. C'est
pour très bientôt.
\end{frame}

